\documentclass[11pt,a4paper]{article}

\usepackage[T1]{fontenc}
\usepackage[utf8]{inputenc}
\usepackage{lmodern}
\usepackage[margin=2.5cm]{geometry}
\usepackage{amsmath,amssymb,amsthm,mathtools}
\usepackage{bm}
\usepackage{booktabs}
\usepackage{array}
\usepackage{xcolor}
\usepackage{listings}
\usepackage{enumitem}
\usepackage[colorlinks=true,linkcolor=blue!50!black,citecolor=blue!50!black,urlcolor=blue!50!black]{hyperref}

\lstdefinestyle{cpp}{
  language=C++,
  basicstyle=\ttfamily\footnotesize,
  keywordstyle=\color{blue!60!black}\bfseries,
  commentstyle=\color{green!40!black}\itshape,
  stringstyle=\color{red!50!black},
  showstringspaces=false,
  breaklines=true,
  frame=single,
  framerule=0.3pt,
  rulecolor=\color{black!30},
  backgroundcolor=\color{black!3},
  xleftmargin=1em,
  framexleftmargin=0.5em,
  columns=fullflexible,
  keepspaces=true
}
\lstset{style=cpp}
\setlength{\emergencystretch}{3em}

\newtheorem{proposition}{Proposition}
\theoremstyle{remark}
\newtheorem*{remark}{Remark}

\newcommand{\R}{\mathbb{R}}
\newcommand{\x}{\bm{x}}
\newcommand{\z}{\bm{z}}
\newcommand{\g}{\bm{g}}
\newcommand{\s}{\bm{s}}
\newcommand{\y}{\bm{y}}
\newcommand{\dd}{\bm{d}}
\newcommand{\Cred}{C_{\mathrm{red}}}
\newcommand{\T}{^{\mathsf T}}
\newcommand{\mT}{^{-\mathsf T}}
\newcommand{\tr}{\operatorname{tr}}
\newcommand{\code}[1]{\texttt{#1}}

\title{Preconditioning L-BFGS with the Analytic FEM Stiffness\\[4pt]
\large Theory and implementation notes for the MTM lattice-triangulation code}
\author{Technical note}
\date{September 2026}

\begin{document}
\maketitle

\begin{abstract}
The quasi-static shear simulations of the MTM code spend almost all of their time
(97\% in a $150\times150$ production run) minimizing the energy on the \emph{elastic
branch}, i.e.\ in load steps where nothing dramatic happens. Plain L-BFGS needs
$\sim10^3$ iterations per such step, and this number grows with the system size.
This note explains, from the ground up, why that happens and how we fixed it: we
build the exact analytic FEM stiffness matrix $K$ from the second derivatives that are
already coded in the energy functions, factorize it with a sparse Cholesky
decomposition once per load step, and feed it to ALGLIB's L-BFGS through a linear
change of variables. ALGLIB itself is not modified. On elastic steps the number of
iterations drops from $\sim1250$ to $\sim37$ at $150\times150$, and the wall time per
step from $5.5$\,s to $0.37$\,s (about $15\times$). Along the way we found and fixed a
factor-of-two inconsistency between the energy and its gradient. Everything is
opt-in (\code{-{}-precond=stiffness}); the default behaviour of the code is unchanged.
The same factorization machinery also computes the lowest eigenvalues of $K$ (stability
analysis) $9$--$60\times$ faster than before (Section~\ref{sec:eig}). Section~\ref{sec:literature}
and the annotated bibliography place all this in the literature.
\end{abstract}

\tableofcontents
\newpage

%=====================================================================
\section{The problem we are trying to solve}
%=====================================================================

Let me start with the physics workflow, because the whole story is about \emph{where}
the time goes.

A simulation (e.g.\ \code{example\_\allowbreak{}1\_\allowbreak{}conti\_\allowbreak{}zanzotto\_\allowbreak{}loading}) applies simple shear
\begin{equation}
F_{\mathrm{ext}}(\alpha) = \begin{pmatrix} 1 & \alpha \\ 0 & 1\end{pmatrix},
\qquad \alpha = 0.14,\ 0.14006,\ 0.14012,\ \dots,\ 1.0,
\end{equation}
in about $14\,334$ small increments $\Delta\alpha = 6\times10^{-5}$. At every load step the
nodes are first mapped affinely by the increment, and then the total energy is
\emph{minimized} over the node positions with L-BFGS (ALGLIB \code{minlbfgs}). This is
the classical quasi-static protocol: the system is assumed to follow a local minimum of
the energy, and when that minimum disappears (an instability) the minimizer carries the
system to a new one. Those events are the avalanches / plastic events we are interested
in.

The key observation is statistical: \emph{avalanches are rare}. In the reference
$150\times150$ run (\code{run\_\allowbreak{}zanzotto\_\allowbreak{}loading\_\allowbreak{}150x150}, seed 42):

\begin{center}
\begin{tabular}{lrr}
\toprule
& steps & time in the main minimization \\
\midrule
elastic steps (no remeshing) & 14\,153 (98.7\%) & 11.3\,h \\
remeshing / avalanche steps  & 180 (1.3\%)     & 0.4\,h \\
\bottomrule
\end{tabular}
\end{center}

So 97\% of the 11.7\,h wall time is spent solving \emph{boring} problems: a lattice that
is already in equilibrium receives a tiny homogeneous shear and has to relax a tiny bit.
Each of those relaxations took on average $\sim2.9$\,s and $\sim10^3$ L-BFGS iterations.
That is absurdly expensive for what is essentially a linear-elastic response, and it is
precisely the regime where second-order information (the Hessian) should help
enormously.

%=====================================================================
\section{The discrete model}
%=====================================================================

\subsection{Kinematics}

The reference lattice is triangulated into P1 (linear) triangles. For an element $e$
with nodes $a=0,1,2$, current positions $\x_a\in\R^2$ and constant shape function
gradients $\nabla N_a = (\partial N_a/\partial X_1,\partial N_a/\partial X_2)$, the
deformation gradient is constant on the element:
\begin{equation}
F = \sum_{a=0}^{2} \x_a \otimes \nabla N_a
\qquad\Longleftrightarrow\qquad
F_{iK} = \sum_a x_{a,i}\,\frac{\partial N_a}{\partial X_K}.
\label{eq:F}
\end{equation}
In the code this is \code{F = x\_\allowbreak{}e * dNdX} with \code{x\_\allowbreak{}e} the $2\times3$ matrix of node
positions and \code{dNdX} the $3\times2$ matrix of shape-function gradients
(\code{ElementTriangle2D::\allowbreak{}calculate\_\allowbreak{}deformation\_\allowbreak{}gradient}). With periodic boundary
conditions a node of an element may be a periodic image; its position is
$\x_{\text{orig}} + F_{\mathrm{ext}}\bm{t}$ with a fixed lattice translation $\bm{t}$.
The translation does not depend on the unknowns, so for all derivatives below a periodic
image behaves exactly like its original node.

The right Cauchy--Green (metric) tensor is $C=F\T F$.

\subsection{Energy and Lagrange reduction}

The strain energy density is the Conti--Zanzotto type function $\varphi(C)$ of the metric,
invariant under the lattice symmetry group $GL(2,\mathbb{Z})$. It is evaluated on the
\emph{reduced} metric
\begin{equation}
\Cred = Z\T C Z = (FZ)\T(FZ), \qquad Z\in GL(2,\mathbb{Z}),
\end{equation}
where $Z$ (\code{m\_\allowbreak{}matrix}) is produced by \code{lagrange::\allowbreak{}reduce(C)} and brings $C$ into
the fundamental domain. $Z$ is piecewise constant: it only jumps when $C$ crosses the
boundary of the fundamental domain, and because $\varphi$ is $GL(2,\mathbb{Z})$-invariant
the energy is continuous across such a jump. Inside each cell of the reduction we can
therefore differentiate with $Z$ held fixed. This is what the gradient code does, and what
the Hessian code must do too.

The total energy is
\begin{equation}
E(\x) = \sum_{e} A_e\,\frac{\varphi\big(\Cred^{(e)}\big) - \varphi_0}{a^2},
\label{eq:E}
\end{equation}
with $A_e$ the reference area (\code{getReferenceArea()}), $a$ the lattice parameter and
$\varphi_0$ the energy of the undeformed crystal.

\subsection{Unknowns}

The free nodes are numbered $0,\dots,n-1$ through \code{full\_\allowbreak{}mapping[node].second}, and
the solver vector is
\begin{equation}
\x = (u_0,\dots,u_{n-1},\ v_0,\dots,v_{n-1}) \in \R^{2n},
\end{equation}
i.e.\ first all $x$-components, then all $y$-components. At $150\times150$ with PBC,
$n=22\,500$ and there are $45\,000$ unknowns; there are $45\,000$ elements.

%=====================================================================
\section{First derivatives: stress and forces}
%=====================================================================

\subsection{The tensor derivative convention}

$\varphi$ is coded as a function of three scalars $(c_{11},c_{22},c_{12})$. A symmetric
matrix has four entries but only three independent ones, so we have to be careful what
``$\partial\varphi/\partial C$'' means. We define the tensor derivative $D$ by requiring
\begin{equation}
d\varphi = \sum_{a,b} D_{ab}\, dC_{ab}\quad\text{for all symmetric } dC.
\end{equation}
Writing it out, $d\varphi = \varphi_{,11}\,dc_{11} + \varphi_{,22}\,dc_{22} + \varphi_{,12}\,dc_{12}$
while $\sum D_{ab}dC_{ab} = D_{11}dc_{11}+D_{22}dc_{22}+(D_{12}+D_{21})dc_{12}$, hence
\begin{equation}
D_{11}=\varphi_{,11},\qquad D_{22}=\varphi_{,22},\qquad D_{12}=D_{21}=\tfrac12\,\varphi_{,12}.
\label{eq:D}
\end{equation}
This is exactly what \code{dphi\_\allowbreak{}func} returns (it computes $\varphi_{,12}$ and then divides
the off-diagonal entry by two: ``we go from DE/DC to J''). I checked it against finite
differences of \code{phi\_\allowbreak{}func}: all three entries agree to $10^{-9}$.

\subsection{First Piola--Kirchhoff stress}

Set $G = FZ$ so that $\Cred = G\T G$, i.e.\ $(\Cred)_{ab} = G_{ia}G_{ib}$ (summation over
repeated indices). Since $G_{ia}=F_{iM}Z_{Ma}$, we have
$\partial G_{ia}/\partial F_{jK} = \delta_{ij}Z_{Ka}$ and
\begin{equation}
B_{ab,jK} \coloneqq \frac{\partial (\Cred)_{ab}}{\partial F_{jK}}
= Z_{Ka}\,G_{jb} + Z_{Kb}\,G_{ja}.
\label{eq:B}
\end{equation}
The stress is then (per unit reference area, dropping $1/a^2$ for readability)
\begin{equation}
P_{jK} = \frac{\partial\varphi}{\partial F_{jK}} = D_{ab}B_{ab,jK}
= (G D\T Z\T)_{jK} + (G D Z\T)_{jK} = 2\,(F Z D Z\T)_{jK},
\label{eq:P}
\end{equation}
which is the line in \code{minimize\_\allowbreak{}energy\_\allowbreak{}with\_\allowbreak{}triangles}:
\begin{lstlisting}
const Eigen::Matrix2d P = 2.0 * F * result.m_matrix * dE_dC
                        * result.m_matrix.transpose();
\end{lstlisting}
It is also useful to name the second Piola--Kirchhoff stress
\begin{equation}
S = 2\,Z D Z\T,\qquad P = FS.
\end{equation}

\subsection{Nodal forces}

From \eqref{eq:F}, $\partial F_{iK}/\partial x_{a,j} = \delta_{ij}\,\partial N_a/\partial X_K$, so
the contribution of element $e$ to the gradient is
\begin{equation}
\frac{\partial E_e}{\partial \x_a} = A_e\, P\,\nabla N_a\T .
\label{eq:force}
\end{equation}
This is \code{ElementTriangle2D::\allowbreak{}calculate\_\allowbreak{}nodal\_\allowbreak{}forces}: \code{(P * dNdX.row(a).transpose()) * area}.

\subsection{A factor of two we found on the way}
\label{sec:area}

Until this work, the line computing $P$ in \code{minimize\_\allowbreak{}energy\_\allowbreak{}with\_\allowbreak{}triangles} ended
with \code{* area}, and \code{calculate\_\allowbreak{}nodal\_\allowbreak{}forces} multiplied by \code{area} again. So
the code returned $A_e\,\partial E/\partial\x$ instead of $\partial E/\partial\x$.

Why did nobody notice? Because on the square lattice every reference triangle has area
exactly $A_e=\tfrac12$. (This is even true after remeshing: a Delaunay triangle of the
lattice contains no other lattice points, and by Pick's theorem such a lattice triangle
has area $\tfrac12$.) So the gradient was exactly $\tfrac12\nabla E$ everywhere: the zeros
of the gradient, i.e.\ the equilibria, were correct, and all the physics was fine.
What was \emph{not} fine is that L-BFGS's line search uses $E$ and $\g$ together (sufficient
decrease and curvature conditions, cubic interpolation), and they disagreed by a factor 2.

We discovered this while validating the Hessian: the directional check
$K\bm{v}\approx[\g(\x+\epsilon\bm{v})-\g(\x-\epsilon\bm{v})]/2\epsilon$ came out as
$0.500000$ for every random direction, and an energy check $dE/dv \approx 2\,\g\cdot\bm{v}$
told us which side was wrong. After removing the extra area (only in
\code{minimize\_\allowbreak{}energy\_\allowbreak{}with\_\allowbreak{}triangles}; the \code{\_\allowbreak{}noreduction} variant was left alone on
purpose) both checks agree to $10^{-7}$, the trajectory of the default solver is unchanged
to $5\times10^{-10}$ in energy, and the plain solver itself became about 16\% faster (fewer
line-search evaluations).

%=====================================================================
\section{Why plain L-BFGS is slow here}
%=====================================================================
\label{sec:slow}

\subsection{L-BFGS in one page}

L-BFGS is a quasi-Newton method. At iteration $k$ it moves along
\begin{equation}
\dd_k = -H_k\,\g_k,
\end{equation}
where $H_k\approx(\nabla^2E)^{-1}$ is an approximation of the inverse Hessian built from
the last $m$ pairs (ALGLIB with $m=13$ in the code)
\begin{equation}
\s_i = \x_{i+1}-\x_i,\qquad \y_i = \g_{i+1}-\g_i .
\end{equation}
$H_k$ is never formed; the product $H_k\g_k$ is computed with the famous
\emph{two-loop recursion}:
\begin{center}
\begin{minipage}{0.8\textwidth}
\begin{tabbing}
\quad\=\quad\=\kill
$\bm q \leftarrow \g_k$\\
for $i=k-1,\dots,k-m$:\> \> $\theta_i = \rho_i\,\s_i\T\bm q,\quad \bm q \leftarrow \bm q - \theta_i\y_i$\qquad ($\rho_i = 1/\y_i\T\s_i$)\\
$\bm r \leftarrow H_0\,\bm q$ \hspace{3em} $\longleftarrow$ \emph{the preconditioner lives here}\\
for $i=k-m,\dots,k-1$:\> \> $\beta = \rho_i\,\y_i\T\bm r,\quad \bm r \leftarrow \bm r + (\theta_i-\beta)\s_i$\\
$\dd_k = -\bm r$
\end{tabbing}
\end{minipage}
\end{center}
The \emph{initial} inverse Hessian $H_0$ is the only free choice. By default ALGLIB uses a
multiple of the identity,
\begin{equation}
H_0 = \gamma_k I,\qquad \gamma_k = \frac{\s_{k-1}\T\y_{k-1}}{\y_{k-1}\T\y_{k-1}},
\end{equation}
which only fixes the overall scale. The step length is then determined by a Mor\'e--Thuente
line search.

\subsection{Condition number of an elastic body}

Near a stable equilibrium the energy is essentially quadratic,
$E\approx\tfrac12(\x-\x^*)\T K(\x-\x^*)$, where $K=\nabla^2E$ is the FEM stiffness matrix.
For linear elasticity discretized with P1 elements on an $L\times L$ grid, $K$ behaves like
a discrete Laplacian. For the 5-point Laplacian with periodic boundary conditions the
eigenvalues are
\begin{equation}
\lambda(k_1,k_2) = 4\sin^2\!\frac{\pi k_1}{L} + 4\sin^2\!\frac{\pi k_2}{L},
\qquad k_1,k_2 = 0,\dots,L-1,
\end{equation}
so (ignoring the zero mode, a rigid translation)
\begin{equation}
\lambda_{\max}\approx 8,\qquad \lambda_{\min}\approx 4\sin^2\frac{\pi}{L}\approx\frac{4\pi^2}{L^2},
\qquad
\kappa(K) = \frac{\lambda_{\max}}{\lambda_{\min}} \approx \frac{2L^2}{\pi^2}.
\end{equation}
For $L=150$ this is $\kappa\approx4.6\times10^3$ for a homogeneous material; the real lattice
with heterogeneous (partly reduced, partly soft) elements is worse.

Gradient-type methods (steepest descent, CG, L-BFGS with a diagonal $H_0$) need a number of
iterations that grows like $\sqrt{\kappa}\propto L$ to damp the smoothest, longest-wavelength
mode. Physically: information travels one element per iteration, and a global shear has to
propagate across the whole box. That is exactly what we see in practice:

\begin{center}
\begin{tabular}{lcc}
\toprule
& $50\times50$ & $150\times150$\\
\midrule
plain L-BFGS, iterations per elastic step & 317 & 1254 \\
\bottomrule
\end{tabular}
\end{center}

A factor 3 in $L$ gives roughly a factor 4 in iterations. The cure is to give L-BFGS a
starting inverse Hessian $H_0$ that already knows about the long-wavelength modes.

\subsection{What ALGLIB offers, and why it is not enough}

\code{minlbfgs} in ALGLIB 3.19 has four modes for $H_0$
(\code{minlbfgssetprec*}):

\begin{center}
\begin{tabular}{p{3.3cm}p{4.6cm}p{5.3cm}}
\toprule
mode & $H_0$ & verdict for us\\
\midrule
default & $\gamma_kI$ & what we used so far\\
diagonal (\code{setprecdiag}) & $D^{-1}$, $D\approx\operatorname{diag}\nabla^2E$ & almost useless (see below)\\
dense Cholesky (\code{setpreccholesky}) & $(U\T U)^{-1}$ with a dense $N\times N$ factor & impossible: $45\,000^2$ doubles $=16$\,GB, $O(N^2)$ work per iteration\\
scale (\code{setprecscale}) & $\operatorname{diag}(s_i^2)$ & useless: all our variables have the same scale\\
\bottomrule
\end{tabular}
\end{center}

\paragraph{Why the diagonal does not help.} With periodic boundary conditions there are no
boundary nodes, and in a homogeneous state every node sees exactly the same patch of
elements. Hence $\operatorname{diag}K$ takes only two values, one for the $u$-components and
one for the $v$-components. A diagonal preconditioner can therefore do nothing but rescale
$x$ against $y$; it has no idea about wavelengths. Worse, looking at the ALGLIB source
(\code{prectype==2} in the two-loop recursion), the diagonal mode \emph{replaces} the
adaptive $\gamma_k$ scaling by the fixed $D^{-1}$. Measured on elastic steps at the same
gradient tolerance, diagonal preconditioning gives $1.3\times$, i.e.\ nothing beyond what
tightening the tolerance does anyway.

\paragraph{A bug in the old helper.} The existing helper
\code{computeDiagonalPreconditioner} (in \code{FEMHessianAssembler}) returned
$1/\sqrt{|H_{ii}|}$. ALGLIB wants the diagonal of the Hessian itself ($H_0=D^{-1}$ is
computed internally), so this would have produced $H_0=\operatorname{diag}\sqrt{H_{ii}}$:
big steps in stiff directions. It is now fixed to return $\max(H_{ii},\text{floor})$, but as
argued above it would not have helped much anyway.

So ALGLIB has no \emph{sparse} preconditioner. We could write our own L-BFGS, but there is a
much cheaper trick.

%=====================================================================
\section{The key idea: preconditioning by a change of variables}
%=====================================================================
\label{sec:change}

\subsection{The transformation}

Let $P$ be any symmetric positive definite matrix that approximates the Hessian, and let
\begin{equation}
P = R\,R\T
\end{equation}
be a factorization (Cholesky). Instead of minimizing $E(\x)$, minimize
\begin{equation}
\tilde E(\z) = E\big(\x_0 + R\mT\z\big),\qquad \z\in\R^{2n},
\label{eq:cov}
\end{equation}
starting from $\z=0$, i.e.\ from $\x=\x_0$. By the chain rule,
\begin{align}
\nabla_{\z}\tilde E &= R^{-1}\,\nabla_{\x}E, \label{eq:gz}\\
\nabla^2_{\z}\tilde E &= R^{-1}\,\nabla^2_{\x}E\;R\mT . \label{eq:Hz}
\end{align}
If $P$ is a good approximation of $\nabla^2E$, then $\nabla^2_{\z}\tilde E\approx
R^{-1}RR\T R\mT = I$: in the $\z$ variables \emph{all directions look equally stiff}. The
long-wavelength modes that were $10^3$--$10^4$ times softer than the short ones are now of
order one. The condition number collapses from $O(L^2)$ to $O(1)$, and L-BFGS on
$\tilde E$ converges in a handful of iterations.

The beauty is that ALGLIB does not need to know anything about this: it just minimizes a
function of $\z$ whose value and gradient we provide in the callback.

\subsection{This is exactly preconditioned L-BFGS}

It is worth convincing ourselves that this is not ``some other method'' but precisely
L-BFGS with $H_0 = \gamma_k P^{-1}$ in the original variables.

\begin{proposition}
Run L-BFGS on $\tilde E(\z)$ with $\tilde H_0 = \gamma I$. Map every iterate back with
$\x = \x_0+R\mT\z$. The resulting sequence $\x_k$ is identical to that of L-BFGS on $E(\x)$
with $H_0 = \gamma P^{-1}$, including the line search.
\end{proposition}

\begin{proof}
Differences transform as $\tilde\s = R\T\s$ and, by \eqref{eq:gz}, $\tilde\y = R^{-1}\y$.
Hence $\tilde\y\T\tilde\s = \y\T R\mT R\T\s = \y\T\s$, so all $\rho_i$ are the same. The
BFGS inverse update in $\z$-space reads
\[
\tilde H_{+} = (I-\rho\,\tilde\s\tilde\y\T)\,\tilde H\,(I-\rho\,\tilde\y\tilde\s\T) + \rho\,\tilde\s\tilde\s\T.
\]
Assume $\tilde H = R\T H R$. Using $I-\rho\tilde\s\tilde\y\T = R\T(I-\rho\s\y\T)R\mT$ and
$I-\rho\tilde\y\tilde\s\T = R^{-1}(I-\rho\y\s\T)R$,
\[
\tilde H_{+} = R\T\Big[(I-\rho\s\y\T)H(I-\rho\y\s\T)+\rho\s\s\T\Big]R = R\T H_{+}R .
\]
By induction $\tilde H_k = R\T H_k R$ for all $k$, provided it holds for $k=0$:
$\tilde H_0=\gamma I = R\T(\gamma R\mT R^{-1})R = R\T(\gamma P^{-1})R$. The search directions
match: $\tilde\dd = -\tilde H\tilde\g = -R\T H R R^{-1}\g = R\T\dd$, i.e.\ $R\mT\tilde\dd = \dd$,
so the trial points along the line are the same in $\x$-space. The function values are the same by
construction, and the directional derivatives too: $\tilde\g\T\tilde\dd = \g\T R\mT R\T\dd =
\g\T\dd$. So every test of the line search gives the same answer.
\end{proof}

Two small remarks. (i) ALGLIB still computes its adaptive scale $\gamma_k$, which in
$\z$-space equals $\s\T\y/(\y\T P^{-1}\y)$. This corrects the overall scale of $P$, so we do
not need to worry about normalizing $P$. (ii) The \emph{stopping criteria} of ALGLIB are
evaluated in $\z$-space: $\|\tilde\g\| = \sqrt{\g\T P^{-1}\g}$ and
$\|\tilde\s\| = \sqrt{\s\T P\s}$, which are energy norms rather than Euclidean norms. That is
why we use our own stopping test in $\x$-space (Section~\ref{sec:stop}).

This idea is not new in atomistic simulation: Packwood, Kermode, Mones, Bernstein, Woolley,
Gould, Ortner and Csányi, \emph{J.\ Chem.\ Phys.}\ \textbf{144}, 164109 (2016), used exactly
such preconditioners (sparse, Laplacian-like, or built from the Hessian) and showed that the
iteration count becomes essentially independent of system size.

\subsection{Which $P$?}

The best $P$ is the Hessian itself at the current configuration, $K=\nabla^2E(\x_0)$. Near a
stable minimum and for a small load increment, the energy is nearly quadratic with Hessian
$K$, so L-BFGS with $H_0=K^{-1}$ is essentially Newton's method (its very first step is the
Newton step) and converges in very few iterations. We compute $K$ exactly and analytically
(Section~\ref{sec:tangent}), with two modifications to make it a valid preconditioner
(Section~\ref{sec:psd}). A cheaper alternative, the reference Laplacian, is discussed in
Section~\ref{sec:laplacian}.

%=====================================================================
\section{The analytic stiffness matrix}
\label{sec:tangent}
%=====================================================================

\subsection{Second derivatives of $\varphi$ with respect to the metric}

The energy file already contains, generated by computer algebra, the six independent
second partials of $\varphi(c_{11},c_{22},c_{12})$:
\[
\varphi_{,1111},\ \varphi_{,2222},\ \varphi_{,1122},\ \varphi_{,1112},\ \varphi_{,2212},\ \varphi_{,1212}
\]
(\code{c1111}, \dots, \code{c1212} in \code{create\_\allowbreak{}d2phi\_\allowbreak{}func}). As for the first derivative,
the tensor second derivative $H_{ab,cd} = \partial D_{ab}/\partial C_{cd}$ is obtained by
differentiating \eqref{eq:D} once more with the same convention. For example
$dD_{11} = \varphi_{,1111}dc_{11}+\varphi_{,1122}dc_{22}+\varphi_{,1112}dc_{12}$, and $dc_{12}$ is
shared between $dC_{12}$ and $dC_{21}$, so $H_{11,12}=H_{11,21}=\tfrac12\varphi_{,1112}$.
Likewise $D_{12}=\tfrac12\varphi_{,12}$ gives $H_{12,12}=H_{12,21}=\tfrac14\varphi_{,1212}$. With
the ordering $(11,12,21,22)$ of the index pairs:
\begin{equation}
H = \begin{pmatrix}
\varphi_{,1111} & \tfrac12\varphi_{,1112} & \tfrac12\varphi_{,1112} & \varphi_{,1122}\\[2pt]
\tfrac12\varphi_{,1112} & \tfrac14\varphi_{,1212} & \tfrac14\varphi_{,1212} & \tfrac12\varphi_{,2212}\\[2pt]
\tfrac12\varphi_{,1112} & \tfrac14\varphi_{,1212} & \tfrac14\varphi_{,1212} & \tfrac12\varphi_{,2212}\\[2pt]
\varphi_{,1122} & \tfrac12\varphi_{,2212} & \tfrac12\varphi_{,2212} & \varphi_{,2222}
\end{pmatrix}.
\label{eq:H}
\end{equation}
This is \emph{exactly} the rule implemented in your \code{AcousticTensor::\allowbreak{}computeHessian}
(``factor 0.5 when one index pair is off-diagonal, 0.25 when both are''), so the physics is
the same as in your existing code. I verified numerically that the coded
\code{ddphi\_\allowbreak{}func} components equal the finite-difference derivatives of \code{dphi\_\allowbreak{}func}
(ratio $1.000000$ for all six). (Finite differences were only used as a \emph{check}; the
tangent itself is fully analytic.)

\subsection{The Lagrangian tangent}

We need $\mathbb{A}_{iK,jL} = \partial^2\varphi/\partial F_{iK}\partial F_{jL} = \partial P_{iK}/\partial F_{jL}$.
Differentiate $P_{iK} = D_{ab}B_{ab,iK}$:
\begin{equation}
\mathbb{A}_{iK,jL} = B_{ab,iK}\,H_{ab,cd}\,B_{cd,jL} \;+\; D_{ab}\,\frac{\partial^2(\Cred)_{ab}}{\partial F_{iK}\partial F_{jL}} .
\end{equation}
From $(\Cred)_{ab}=G_{ma}G_{mb}$ and $\partial G_{ma}/\partial F_{iK} = \delta_{mi}Z_{Ka}$,
\begin{align*}
\frac{\partial^2(\Cred)_{ab}}{\partial F_{iK}\partial F_{jL}} &= \delta_{ij}\,(Z_{Ka}Z_{Lb}+Z_{Kb}Z_{La}),\\
D_{ab}\,\delta_{ij}(Z_{Ka}Z_{Lb}+Z_{Kb}Z_{La}) &= \delta_{ij}\,2(ZDZ\T)_{KL} = \delta_{ij}S_{KL}.
\end{align*}
So the tangent is the sum of a \emph{material} part and a \emph{geometric} (initial-stress)
part:
\begin{equation}
\boxed{\;\mathbb{A} = B\T H B \;+\; I\otimes S,\qquad
B_{ab,iK} = Z_{Ka}G_{ib}+Z_{Kb}G_{ia},\quad G=FZ,\quad S=2ZDZ\T.\;}
\label{eq:A}
\end{equation}
Here $\mathbb{A}$ is stored as a $4\times4$ matrix with row index $2i+K$ and column index
$2j+L$, $B$ as a $4\times4$ matrix with row $2a+b$ and column $2i+K$, and $H$ as the
$4\times4$ matrix \eqref{eq:H} divided by $a^2$.

In your \code{AcousticTensor} this same object is built with ITensor contractions: the
tensor \code{T53} is $\partial C_{mn}/\partial F_{ij} = \delta_{mj}F_{in}+\delta_{nj}F_{im}$ (contracted
with $Z$ it becomes our $B$), \code{T54} produces the $\delta_{ij}S_{KL}$ term, \code{mm1\dots mm4}
are the $Z$ matrices, and \code{f1+f2} is the sum in \eqref{eq:A}. The new code computes
\eqref{eq:A} with fixed-size Eigen matrices:
\begin{lstlisting}
Eigen::Matrix4d element_lagrangian_tangent(const Eigen::Matrix2d &F,
    const Eigen::Matrix2d &Z, const Eigen::Matrix2d &dE_dC_reduced,
    const HessianComponents &h, double normalisation) {
  // H: tensor second derivative, same 1/2, 1/4 factors as AcousticTensor::computeHessian
  const double d1112 = 0.5 * h.c11_c12, d2212 = 0.5 * h.c22_c12, d1212 = 0.25 * h.c12_c12;
  Eigen::Matrix4d H;
  H << h.c11_c11, d1112, d1112, h.c11_c22,
       d1112,     d1212, d1212, d2212,
       d1112,     d1212, d1212, d2212,
       h.c11_c22, d2212, d2212, h.c22_c22;
  H /= normalisation;
  // B(2a+b, 2i+K) = d(C_red)_ab / dF_iK = Z_Ka G_ib + Z_Kb G_ia,  G = F Z
  const Eigen::Matrix2d G = F * Z;
  Eigen::Matrix4d B;
  for (int a = 0; a < 2; a++) for (int b = 0; b < 2; b++)
    for (int i = 0; i < 2; i++) for (int K = 0; K < 2; K++)
      B(2*a+b, 2*i+K) = Z(K,a)*G(i,b) + Z(K,b)*G(i,a);
  Eigen::Matrix4d A = B.transpose() * H * B;
  // geometric term delta_ij S_KL, S = 2 Z (dE/dC) Z^T
  const Eigen::Matrix2d S = 2.0 * Z * dE_dC_reduced * Z.transpose();
  for (int i = 0; i < 2; i++) for (int K = 0; K < 2; K++) for (int L = 0; L < 2; L++)
    A(2*i+K, 2*i+L) += S(K,L);
  return A;
}
\end{lstlisting}

\subsection{From the tangent to the element stiffness}

Using $\partial F_{iK}/\partial x_{a,j} = \delta_{ij}\,\partial N_a/\partial X_K$ twice,
\begin{equation}
K^{e}_{(a,i),(b,j)} = \frac{\partial^2 E_e}{\partial x_{a,i}\partial x_{b,j}}
= A_e\sum_{K,L}\mathbb{A}_{iK,jL}\,\frac{\partial N_a}{\partial X_K}\,\frac{\partial N_b}{\partial X_L}.
\end{equation}
In matrix form, with the $4\times6$ matrix $\mathcal{G}_{2i+K,\,2a+j} = \delta_{ij}\,\partial N_a/\partial X_K$,
\begin{equation}
K^{e} = A_e\,\mathcal{G}\T\,\mathbb{A}\,\mathcal{G}\qquad (6\times6).
\label{eq:Ke}
\end{equation}
This is the same formula as in \code{FEMHessianAssembler::\allowbreak{}computeElementStiffness}.

\subsection{Getting the second derivatives without ITensor}

The existing \code{BaseLatticeCalculator::\allowbreak{}calculate\_\allowbreak{}dseconderivative} returns an ITensor. Building
and reading ITensors per element is what made the old assembler slow (7.2\,s serially for
5000 elements). So I added one virtual method,
\begin{lstlisting}
virtual HessianComponents calculate_dseconderivative_components(
    const Eigen::Matrix2d& C, const std::function<double(double)>& dpot,
    const std::function<double(double)>& d2pot);
\end{lstlisting}
which returns the six raw partials in the \code{HessianComponents} struct that already existed
in the base header. The base-class default reads them back from the ITensor (so any future
calculator works automatically); \code{Strain\_\allowbreak{}Energy\_\allowbreak{}LatticeCalculator},
\code{SquareLatticeCalculator} and \code{TriangularLatticeCalculator} override it and never touch
ITensor. For the strain energy, the huge generated expressions were moved into
\code{create\_\allowbreak{}d2phi\_\allowbreak{}components\_\allowbreak{}func}; \code{create\_\allowbreak{}d2phi\_\allowbreak{}func} now calls it and packs the ITensor,
so the expressions exist only once.

\subsection{Global assembly}

$K$ is the $2n\times2n$ sparse matrix $\sum_e (\text{scatter of }K^e)$. The DOF of local entry
$(a,i)$ is $s_a$ for $i=0$ and $s_a+n$ for $i=1$, with $s_a=$ \code{full\_\allowbreak{}mapping[nn[a]].second}.
Couplings across the periodic boundary come for free because images carry the DOF of their
original node.

The implementation (\code{FastStiffnessAssembler::\allowbreak{}assembleStiffness}) is organized for speed:
\begin{enumerate}[leftmargin=*]
\item \textbf{Pattern once.} The sparsity pattern depends only on the connectivity. It is built
once (with \code{setFromTriplets}) and, for every element, the positions of its 36 entries in
the compressed value array are stored (\code{k\_\allowbreak{}slots\_\allowbreak{}}). A hash of the connectivity
(\code{mesh\_\allowbreak{}signature}) detects when remeshing changes the mesh, and only then is the pattern
rebuilt.
\item \textbf{Values in parallel.} An OpenMP loop over elements computes $F$, reduces $C$,
evaluates $D$ and $H$, forms $\mathbb{A}$ and $K^e$, and adds the 36 numbers directly into the
value array with \code{\#pragma omp atomic}.
\end{enumerate}
At $50\times50$: 2.9\,ms the first time (pattern included), 0.4\,ms afterwards, versus 7.2\,s for
the serial ITensor loop. At $150\times150$ assembly takes about 8\,ms.

\subsection{Validation of $K$}

The benchmark executable (\code{benchmark\_\allowbreak{}preconditioner -{}-validate}) checks, on a real relaxed
configuration with 3106 of 5000 elements Lagrange-reduced:
\begin{itemize}[leftmargin=*]
\item element by element against the ITensor \code{computeElementStiffness}:
$\max|K^e_{\text{fast}}-K^e_{\text{ITensor}}| / \max|K^e| = 3\times10^{-16}$;
\item symmetry: $\max|K-K\T| = 10^{-14}$;
\item $K$ against the analytic gradient: $\|K\bm v - [\g(\x+\epsilon\bm v)-\g(\x-\epsilon\bm v)]/2\epsilon\|/\|K\bm v\| \approx 10^{-7}$
for random unit $\bm v$ (the finite-difference noise level);
\item gradient against the energy, at a non-equilibrium point: $[E(\x+\epsilon\bm v)-E(\x-\epsilon\bm v)]/(2\epsilon\,\g\cdot\bm v) = 1.0000$.
\end{itemize}
In other words, $K$ is the exact Hessian of the discrete energy \eqref{eq:E}, the gradient is its
exact first derivative, and both coincide with your existing ITensor implementation.

%=====================================================================
\section{Turning $K$ into a preconditioner}
%=====================================================================
\label{sec:psd}

Two things prevent us from Cholesky-factorizing the exact Hessian directly.

\subsection{Indefiniteness: element-wise PSD projection}

Near instabilities, and in the strongly non-equilibrium states after noise or during an
avalanche, some elements are in a non-convex part of $\varphi$ and $\mathbb{A}$ has negative
eigenvalues. Then $K$ is indefinite, Cholesky fails, and even if it did not, $H_0=K^{-1}$ would
not be positive definite and the L-BFGS direction might point uphill.

The fix is a standard trick from projected Newton methods: project every \emph{element}
tangent onto the positive semi-definite cone before assembly. $\mathbb{A}$ is symmetric
(major symmetry of a second derivative), so
\begin{equation}
\mathbb{A} = Q\Lambda Q\T \quad\longrightarrow\quad \mathbb{A}^{+} = Q\max(\Lambda,0)\,Q\T .
\end{equation}
Then $K^{e+} = A_e\mathcal{G}\T\mathbb{A}^+\mathcal{G}$ is PSD, and a sum of PSD matrices is PSD, so
the assembled $K^+$ is PSD. Where the material is stable nothing changes; where an element
is unstable its unstable direction simply gets zero stiffness in the preconditioner. That is
actually a sensible thing to tell the optimizer: ``you may take large steps in this
direction''. The line search guarantees descent anyway.

The $4\times4$ symmetric eigen-decomposition per element costs almost nothing
(\code{Eigen::\allowbreak{}SelfAdjointEigenSolver<Matrix4d>}).

\subsection{Rigid translations: a small shift}

With periodic boundary conditions a rigid translation of all nodes does not change any $F$,
hence $K\bm{t}=0$ for the two translation vectors $\bm t=(1,\dots,1,0,\dots,0)$ and
$(0,\dots,0,1,\dots,1)$: $K$ has a two-dimensional null space. We factorize
\begin{equation}
P = K^{+} + \sigma I,\qquad \sigma = 10^{-6}\,\overline{\operatorname{diag}K^{+}} .
\end{equation}
The shift is far below the smallest non-zero eigenvalue ($\sim(2\pi/L)^2$ times the typical
stiffness, i.e.\ $\sim10^{-3}$ relative at $L=150$), so it does not spoil the preconditioner.
It is harmless for the translation directions: the net force $\sum_a \partial E/\partial\x_a$
vanishes identically (because $\sum_a\nabla N_a=0$ on every element), so the gradient has no
component along the null space. If a factorization nevertheless fails, the code multiplies
$\sigma$ by 100 and retries (up to 6 times), and as a last resort falls back to plain L-BFGS.

%=====================================================================
\section{Sparse Cholesky and the two half-solves}
%=====================================================================
\label{sec:cholesky}

\subsection{Factorization}

Sparse Cholesky first chooses a fill-reducing permutation $\Pi$ (AMD ordering, computed once
per mesh in the ``analyze'' phase) and then computes a lower-triangular $L$ with
\begin{equation}
\Pi\,P\,\Pi\T = L L\T .
\end{equation}
Hence $P = R R\T$ with
\begin{equation}
R = \Pi\T L,\qquad R^{-1} = L^{-1}\Pi,\qquad R\mT = \Pi\T L\mT
\end{equation}
(a permutation matrix is orthogonal, $\Pi^{-1}=\Pi\T$).

Two back-ends are implemented behind a small interface (\code{SparseCholeskyBackend}):
\begin{itemize}[leftmargin=*]
\item \textbf{CHOLMOD} (SuiteSparse, supernodal). Columns of $L$ with identical sparsity are
grouped into dense ``supernodes'' that are processed with BLAS-3, which is fast and can be
multithreaded. Used when CMake / \code{Makefile.cluster} finds SuiteSparse
(\code{-DUSE\_\allowbreak{}CHOLMOD=1}).
\item \textbf{Eigen \code{SimplicialLLT}} (column by column, serial). Always available; the
fallback.
\end{itemize}
Both use the shift directly in the numeric factorization (\code{cholmod\_\allowbreak{}factorize\_\allowbreak{}p} with
$\beta=\sigma$, resp.\ \code{SimplicialLLT::\allowbreak{}setShift}), so $K$ itself is never modified.

Measured at $150\times150$ ($45\,000$ unknowns): factorization 0.032\,s with CHOLMOD, 0.33\,s
with Eigen.

\subsection{Applying $R^{-1}$ and $R\mT$}

We never form $R^{-1}$. The two operations needed per energy evaluation are
\begin{align}
\text{gradient to }\z:\quad & \tilde\g = R^{-1}\g = L^{-1}(\Pi\g)
&& \text{(permute, then forward substitution)}\\
\text{step back to }\x:\quad & \x-\x_0 = R\mT\z = \Pi\T(L\mT\z)
&& \text{(back substitution, then inverse permutation)}
\end{align}
With CHOLMOD these are \code{cholmod\_\allowbreak{}solve2} calls with \code{CHOLMOD\_\allowbreak{}P}, \code{CHOLMOD\_\allowbreak{}L}
and \code{CHOLMOD\_\allowbreak{}Lt}, \code{CHOLMOD\_\allowbreak{}Pt}; with Eigen, \code{permutationP()},
\code{matrixL().solveInPlace()} and \code{matrixU().solveInPlace()},
\code{permutationPinv()}. Each triangular solve costs $O(\operatorname{nnz}(L))$; together
about 4\,ms at $150\times150$, versus about 2\,ms for one energy+gradient evaluation.

%=====================================================================
\section{The algorithm as implemented}
%=====================================================================

\subsection{One call of \code{PreconditionedLBFGS::optimize}}

For every relaxation (one load step, or one relaxation inside the remeshing loop):
\begin{enumerate}[leftmargin=*]
\item \textbf{Build the preconditioner} at the starting point $\x_0$ (\code{buildPreconditioner}):
assemble $K^+(\x_0)$ ($\approx8$\,ms), factorize $K^++\sigma I$ ($\approx32$\,ms). The
symbolic analysis (ordering) is redone only when the mesh changes.
\item \textbf{Run ALGLIB on $\z$} starting from $\z=0$, with the callback
\code{transformed\_\allowbreak{}energy}:
\begin{lstlisting}
static void transformed_energy(const alglib::real_1d_array &z, double &func,
                               alglib::real_1d_array &grad, void *ptr) {
  auto *ctx = static_cast<TransformContext *>(ptr);
  double *x = ctx->x_work.getcontent();
  factor->apply_inverse_transpose(z.getcontent(), x, n);   // x - x0 = R^{-T} z
  for (int i = 0; i < n; i++) x[i] += ctx->x0[i];
  minimize_energy_with_triangles(ctx->x_work, func, ctx->gx_work, ctx->userData);
  factor->apply_inverse(ctx->gx_work.getcontent(), grad.getcontent(), n); // R^{-1} g
  // stop on the true force, measured in x-space (see below)
  if (max_i |g_i| < grad_tol) { remember x; alglib::minlbfgsrequesttermination(*ctx->state); }
}
\end{lstlisting}
Your energy routine \code{minimize\_\allowbreak{}energy\_\allowbreak{}with\_\allowbreak{}triangles} is called unchanged.
\item \textbf{Map back}: $\x = \x_0 + R\mT\z_{\text{final}}$ and return it in the same array the
old optimizer used, so all the code after the minimization (remeshing decision, stresses,
output) is untouched.
\end{enumerate}
$K$ is \emph{not} updated during the ALGLIB iterations. What adapts is the L-BFGS memory
(the 13 pairs), which corrects $H_0=K^{-1}$ with the curvature actually encountered. $K$ cannot
be swapped in the middle of a run: a new $R$ would define new variables $\z$ and make the stored
pairs meaningless. Changing $R$ therefore always means restarting ALGLIB, which is what the
restart logic below does.

\subsection{Stopping criterion}
\label{sec:stop}

With the change of variables, ALGLIB's own criteria would be measured in $\z$-space (energy
norms, see above), and its automatic criterion (a small step, $\|\tilde\s\|\le10^{-6}$) would
mean something different from the production runs. Instead we stop when the true force is
small,
\begin{equation}
\max_i\Big|\frac{\partial E}{\partial x_i}\Big| < \texttt{grad\_tol} = 10^{-6},
\end{equation}
checked inside the callback; ALGLIB is then asked to terminate
(\code{minlbfgsrequesttermination}) and the point that met the criterion is returned. The same
test is used for all preconditioners, so comparisons are fair. (For reference, the plain
production runs stop at $\max|g|$ between $10^{-8}$ and $5\times10^{-5}$; energies of the two
solvers agree to $10^{-10}$.)

\subsection{Restarts when the preconditioner goes stale}

During a large rearrangement (an avalanche, or relaxing the initial noise) the configuration
moves far away from $\x_0$ and $K(\x_0)$ becomes a poor model. ALGLIB may then stop early for
numerical reasons without reaching \code{grad\_\allowbreak{}tol}. In that case \code{optimize} rebuilds
$K$ at the current point, refactorizes, and starts a new ALGLIB run from there (at most 3
restarts). In the tests below no restart was ever needed on avalanche steps; it is a safety
net.

\subsection{Refresh policy}

By default ($\texttt{refresh\_every}=1$) $K$ is rebuilt at the start of every load step. One can
reuse the factorization for several steps (\code{-{}-precond-refresh=N}), optionally refreshing
early when a solve needed more than a given number of iterations. With CHOLMOD the setup is so
cheap (40\,ms) that this buys nothing measurable ($14.9\times$ vs $15.2\times$), so the default
is to refresh every step.

\subsection{The Laplacian variant}
\label{sec:laplacian}

A much simpler preconditioner uses only the mesh: the scalar P1 Laplacian of the reference
configuration,
\begin{equation}
\mathcal{L}_{s_a s_b} = \sum_e A_e\,\nabla N_a\cdot\nabla N_b,\qquad
P = \begin{pmatrix}\mathcal{L}+\sigma I & 0\\ 0&\mathcal{L}+\sigma I\end{pmatrix},
\end{equation}
which does not depend on $\x$ at all, is factorized once per mesh, and acts on $u$ and $v$
separately with the same $n\times n$ factor. It captures the long-wavelength problem (the
$L^2$ condition number) but knows nothing about the actual elastic constants, anisotropy or
the reduced elements. It gives about $3\times$; the true stiffness gives $15\times$.

\subsection{Cost of one elastic step at $150\times150$}

\begin{center}
\begin{tabular}{lrr}
\toprule
& plain L-BFGS & stiffness preconditioner \\
\midrule
assembly of $K$ & -- & 0.008\,s \\
factorization (CHOLMOD) & -- & 0.032\,s \\
iterations & $\sim1250$ & $\sim37$ \\
function evaluations & $\sim2490$ & $\sim53$ \\
cost per evaluation & $\approx2.2$\,ms & $\approx6.2$\,ms \\
\midrule
total per step & 5.5\,s & 0.37\,s \\
\bottomrule
\end{tabular}
\end{center}
Each evaluation is about three times more expensive (two triangular solves on top of the
energy), but we need about 40 times fewer of them.

%=====================================================================
\section{Results}
%=====================================================================

All benchmarks were run with \code{benchmark\_\allowbreak{}preconditioner}. At each load step all methods
start from the \emph{same} configuration, and the state is then advanced with the solution of
the first method (plain production L-BFGS), so every method solves exactly the same sequence of
problems. The loading is the production Zanzotto shear, PBC, fixed mesh, M1 Pro with 10 threads.

\subsection{Elastic steps}

\begin{center}
\begin{tabular}{lrrrr}
\toprule
& \multicolumn{2}{c}{$50\times50$} & \multicolumn{2}{c}{$150\times150$ (CHOLMOD)}\\
\cmidrule(lr){2-3}\cmidrule(lr){4-5}
method & its/step & speedup & its/step & speedup\\
\midrule
plain (production) & 245 & 1.0 & 1254 & 1.0\\
plain, same tolerance & 228 & 1.1 & -- & --\\
diagonal (ALGLIB) & 237$^\ast$ & 1.3$^\ast$ & -- & --\\
Laplacian & 64 & 3.2 & 207$^\dagger$ & 3.0$^\dagger$\\
\textbf{stiffness} & \textbf{10.7} & \textbf{12.3} & \textbf{37} & \textbf{14.9}\\
\bottomrule
\end{tabular}
\end{center}
{\footnotesize $^\ast$ measured before the gradient fix; $^\dagger$ with the Eigen back-end run.
With Eigen instead of CHOLMOD the stiffness preconditioner gives $7.8\times$ at $150\times150$, the
difference being entirely the factorization time.}

In all cases the preconditioned solutions coincide with the production solutions to about
$10^{-10}$ in energy. Note that the iteration count of the preconditioned solver still grows
somewhat with size ($11\to37$): $K(\x_0)$ is not the Hessian at the solution, the problem is
nonlinear, and the absolute force tolerance involves more degrees of freedom. But the gap to
plain L-BFGS widens with size, which is what matters for large systems.

\subsection{Avalanches}

To see many instabilities quickly, a $50\times50$ run with ten times larger increments
($\Delta\alpha=6\times10^{-4}$, 150 steps, $\alpha$ from 0.14 to 0.23) was used. It contains 9
energy drops. Total time excluding step 0:

\begin{center}
\begin{tabular}{lrrr}
\toprule
method & elastic steps (140) & avalanche steps (9) & total \\
\midrule
plain & 43.5\,s & 5.8\,s & 49.2\,s\\
Laplacian & 15.2\,s & 2.1\,s & 17.4\,s\\
\textbf{stiffness} & \textbf{2.8\,s} & \textbf{1.1\,s} & \textbf{3.9\,s ($12.6\times$)}\\
\bottomrule
\end{tabular}
\end{center}
On the avalanche steps themselves the stiffness solver needed 50--153 iterations against
650--1180 for plain L-BFGS, with no restarts. In 7 of the 9 avalanches it ended in exactly the
same state as plain L-BFGS ($|\Delta E|\sim10^{-11}$). In two it did not: once it found a state
lower by $0.076$ (the Laplacian found the same one), once a state higher by $0.047$. This is
discussed in Section~\ref{sec:physics}.

\subsection{The production binary}

Running the real program for 45\,s at $50\times50$:
\begin{center}
\begin{tabular}{lr}
\toprule
configuration & load steps in 45\,s\\
\midrule
plain, before the gradient fix & 182\\
plain, after the gradient fix & 212\\
\code{-{}-precond=stiffness -{}-precond-from-step=1} & 2062\\
\bottomrule
\end{tabular}
\end{center}
With \code{-{}-precond-from-step=1} the first 182 steps match the plain run to $1.2\times10^{-10}$
in energy and $1.2\times10^{-7}$ in stress.

A full $150\times150$ production-style run (positive shear, seed 42, remeshing on) was started
with the stiffness preconditioner: after 50 minutes it had completed 6586 steps, with elastic
steps at 0.30\,s against 2.87\,s in the old run ($9.6\times$, including remeshing overhead), on
track for $\sim1.8$\,h instead of 11.7\,h.

%=====================================================================
\section{Physics caveats: path dependence}
\label{sec:physics}
%=====================================================================

A preconditioner changes the \emph{path} of the minimizer, not the energy landscape. On a convex
problem the path is irrelevant: there is one minimum and every method finds it. That is why
elastic steps agree to $10^{-10}$. But whenever the starting point is far from equilibrium in a
non-convex landscape, different paths may end in different local minima. Neither answer is
``wrong'': quasi-static minimization does not define which basin the system falls into after an
instability; that is decided by the (fictitious) minimization dynamics.

This matters in two places.

\paragraph{Step 0.} The first load step relaxes a lattice sheared to $\alpha=0.14$ with 4\%
random noise on the positions: a violent, strongly non-convex relaxation ($\sim5000$ iterations
at $150\times150$). At $50\times50$, plain L-BFGS ends at $E=14.09$ with shear stress
$+7.1\times10^{-3}$, while the preconditioned solver ends at $E=7.99$ with stress
$-4.1\times10^{-3}$: a different initial microstructure, i.e.\ effectively a different sample.
The option \code{-{}-precond-from-step=N} runs the plain solver (including its remeshing) for load
steps $<N$. With $N=1$ (the default in the job scripts) the initial state is identical to that
of a plain run, at a cost of about 20\,s out of hours.

(A side remark: at $150\times150$ the step-0 energy of today's code, with or without
preconditioner, is 122.805, while the old run from 9--10 September gave 116.062. The step-0
relaxation is so sensitive that the code changes made since then, this morning's optimization of
\code{minimize\_\allowbreak{}energy\_\allowbreak{}with\_\allowbreak{}triangles} and \code{lagrange::\allowbreak{}reduce} and/or the gradient fix, send
it to a different minimum. New runs should therefore be compared with new plain runs, or
statistically.)

\paragraph{Avalanches.} As seen above, most avalanches end in the same state, but not all.
The scientific output of these simulations is statistical (avalanche size distributions, stress
drops, dislocation statistics), so the right validation is statistical: run a few seeds with and
without preconditioner and compare the distributions.

%=====================================================================
\section{Stability analysis with the same machinery: the lowest eigenvalues of $K$}
\label{sec:eig}
%=====================================================================

\subsection{What we want and how it was done before}

The \emph{unprojected} stiffness $K$ is the Hessian of the discrete energy. At a relaxed
configuration all its eigenvalues are $\ge0$; an avalanche starts when the lowest nontrivial
eigenvalue $\lambda_{\min}$ crosses zero (a saddle-node bifurcation), and the corresponding
eigenvector is the soft mode that becomes unstable. The analysis code
(\code{analyze\_\allowbreak{}data\_\allowbreak{}from\_\allowbreak{}folder} in \code{src/\allowbreak{}experiments/\allowbreak{}data\_\allowbreak{}analysis.cpp})
therefore loads saved configurations and computes the 100 lowest eigenpairs of $K$:
\begin{enumerate}[leftmargin=*]
\item $K$ is assembled by \code{FEMHessianAssembler::\allowbreak{}assembleGlobalStiffness}, element by element
through ITensor, serially;
\item \code{computeSmallestEigenvaluesIterative\_\allowbreak{}spectra} runs Spectra's shift-invert Lanczos
with \code{SparseSymShiftSolve}, which factorizes $K-\sigma I$ with a sparse \emph{LU}
(\code{Eigen::\allowbreak{}SparseLU}), with $\sigma=-0.01\max_i|K_{ii}|$;
\item the two rigid translations (eigenvalue $0$ under periodic boundary conditions) come out
among the 100 and are recognized afterwards (\code{detectRigidBodyModes}).
\end{enumerate}
At $150\times150$ this takes 20\,s per configuration: 3.7\,s for the ITensor assembly and 16\,s
for the eigensolver (with the release ITensor library \code{libitensor.a}; with the debug library
\code{libitensor-g.a}, which \code{CMakeLists.txt} used to prefer and now links only with
\code{-DITENSOR\_\allowbreak{}DEBUG\_\allowbreak{}LIB=ON}, the assembly alone takes 65\,s).

\subsection{Shift-invert Lanczos in one paragraph}

Lanczos finds the \emph{extreme} eigenvalues of a symmetric operator quickly, but the lowest
eigenvalues of a stiffness matrix are extreme only in a useless sense: they are tiny and
clustered compared with the spread $\lambda_{\max}\sim8\bar K_{ii}$. The spectral
transformation~\cite{ericsson1980} applies Lanczos to $(K-\sigma I)^{-1}$ instead, whose
eigenvalues are
\begin{equation}
\nu_k = \frac{1}{\lambda_k-\sigma},\qquad \lambda_k = \sigma + \frac1{\nu_k}.
\end{equation}
The eigenvalues closest to $\sigma$ become the largest $\nu$'s and are well separated from the
rest, so they converge in a few tens of iterations. Every Lanczos iteration costs one solve with
$K-\sigma I$, so the method lives or dies with the factorization of $K-\sigma I$.

\subsection{Three improvements}

\paragraph{1. Fast assembly.} $K$ is assembled by the \code{FastStiffnessAssembler} of
Section~\ref{sec:tangent} with \code{project\_\allowbreak{}psd=false}: exactly the same matrix (checked:
$\max|K_{\text{fast}}-K_{\text{ITensor}}|=7\times10^{-13}$), in 25\,ms instead of 3.7\,s.

\paragraph{2. Cholesky instead of LU: put the shift \emph{below} the spectrum.} If
$\sigma<\lambda_{\min}(K)$, then $K-\sigma I$ is symmetric positive definite and we can use the
sparse Cholesky factorization of Section~\ref{sec:cholesky} (CHOLMOD, 45\,ms at $150\times150$), which is several
times cheaper than LU in time and memory and gives faster solves. We do not know
$\lambda_{\min}$ in advance, but we do not need to: by Sylvester's law of
inertia~\cite{grimes1994}, the number of negative pivots of a symmetric factorization of
$K-\sigma I$ equals the number of eigenvalues of $K$ below $\sigma$. In particular
\begin{equation}
\text{Cholesky of }K-\sigma I\text{ succeeds}\iff \lambda_{\min}(K)>\sigma .
\end{equation}
So we try $\sigma_0=-10^{-4}\,\overline{K_{ii}}$ and, while the factorization fails, set
$\sigma\leftarrow10\,\sigma$. A failure is not wasted: it is a rigorous statement that $K$ has an
eigenvalue below $\sigma$, i.e.\ \emph{the configuration is unstable}. After $m$ failures,
$\lambda_{\min}\in(\sigma,\sigma/10]$, so the shift is always within a factor 10 of
$|\lambda_{\min}|$ and the lowest modes converge fast.

\paragraph{3. Deflate the translations instead of computing them.} Under PBC every node is free,
and a uniform translation $\bm t_u=(1,\dots,1,0,\dots,0)$ or $\bm t_v=(0,\dots,0,1,\dots,1)$ does
not change any deformation gradient, so $K\bm t_u=K\bm t_v=0$ exactly. They are therefore also
eigenvectors of $(K-\sigma I)^{-1}$, and with the orthogonal projector
$Q=I-\hat{\bm t}_u\hat{\bm t}_u\T-\hat{\bm t}_v\hat{\bm t}_v\T$ (which just subtracts the mean of
the $u$ block and of the $v$ block) the operator
\begin{equation}
\mathcal{A}\,\bm x = Q\,(K-\sigma I)^{-1}\,Q\,\bm x
\end{equation}
is symmetric, has the same eigenpairs as $(K-\sigma I)^{-1}$ on the complement of the
translations, and maps the translations to $0$, where Lanczos never looks. No eigenpairs are
spent on them and no detection step is needed. There is a second benefit: the translation
eigenvalue $0$ is double, and a single-vector Lanczos can return only one copy of a double
eigenvalue. With the old code and 3 requested values this happens: one translation, $\lambda_1$
and $\lambda_2$ come back, so the number of ``rigid modes'' found depends on how many values are
requested. The code checks $\|K\bm t\|_\infty\le10^{-9}\max|K_{ij}|$
before deflating, so with fixed boundaries (no translation modes) nothing is projected.

\subsection{Results}

Same relaxed $150\times150$ configuration ($n=45\,000$; 3 load steps of the benchmark), M1 Pro,
CHOLMOD, release ITensor library for the old path:

\begin{center}
\begin{tabular}{lrrr}
\toprule
values requested & 100 & 12 & 3\\
(nontrivial modes) & (98) & (10) & (1: $\lambda_{\min}$)\\
\midrule
old: ITensor assembly & 3.7\,s & 3.7\,s & 3.6\,s\\
old: Spectra + SparseLU & 16.2\,s & 6.9\,s & 5.2\,s\\
\textbf{old total} & \textbf{19.9\,s} & \textbf{10.5\,s} & \textbf{8.8\,s}\\
\midrule
new: assembly & 0.03\,s & 0.02\,s & 0.02\,s\\
new: Cholesky & 0.05\,s & 0.04\,s & 0.04\,s\\
new: Lanczos (solves) & 2.22\,s (248) & 0.18\,s (45) & 0.08\,s (21)\\
\textbf{new total} & \textbf{2.32\,s} & \textbf{0.25\,s} & \textbf{0.15\,s}\\
\midrule
speedup & $8.6\times$ & $42\times$ & $59\times$\\
\bottomrule
\end{tabular}
\end{center}

The eigenvalues agree with the old ones to a relative $1.5\times10^{-11}$ (98 modes), and the
residuals $\|K\bm v-\lambda\bm v\|/|\lambda|$ are $\le1.5\times10^{-11}$. With the debug ITensor
library the old path takes 79\,s for 100 values. With $\lambda_{\min}$ at 0.15\,s it becomes
affordable to monitor stability \emph{at every load step} of a simulation, not only at a few saved
configurations.

\paragraph{In \code{analyze\_\allowbreak{}data\_\allowbreak{}from\_\allowbreak{}folder}.} On saved configurations 1, 3, 5 of
the $150\times150$ preconditioned run (100 values, release ITensor), the eigenvalue part takes
2.2\,s (fast) against 14--15\,s (legacy) per configuration; the output file
\code{eigenvalues\_\allowbreak{}vs\_\allowbreak{}shear.dat} is the same to 11 digits in the first nontrivial
eigenvalue, with 2 rigid modes detected in both cases.

\paragraph{Unstable test.} To exercise the shift search, the solver was applied to
$K-cI$ with $c$ chosen so that three eigenvalues are negative. The first shift
$\sigma_0=-5.5\times10^{-3}$ fails (instability detected), $\sigma=-5.5\times10^{-2}$ succeeds, and
the five lowest eigenvalues $-3.389\times10^{-2}$, $-1.122\times10^{-3}$, $-8.776\times10^{-4}$,
$1.018\times10^{-2}$, $1.891\times10^{-2}$ are reproduced to all printed digits.

\paragraph{Usage.} From C++:
\begin{lstlisting}
#include "optimization/StiffnessSpectrum.h"

StiffnessSpectrumOptions opt;
opt.n_modes = 10;                 // lowest 10 nontrivial modes
StiffnessSpectrum s = lowest_stiffness_modes(x, &userData, opt);
// s.eigenvalues (ascending), s.eigenvectors (solver layout [u..., v...]),
// s.failed_shifts > 0  <=>  K has an eigenvalue below the first trial shift
\end{lstlisting}
The command \code{benchmark\_\allowbreak{}preconditioner -{}-eig\allowbreak{} 100} reproduces the comparison above.
The function \code{analyze\_\allowbreak{}data\_\allowbreak{}from\_\allowbreak{}folder} uses the fast path by default and the old
one with \code{-{}-eig-\allowbreak{}solver=\allowbreak{}legacy} (flag of \code{lattice\_\allowbreak{}triangulation}, passed as the last
argument of the call in \code{main.cpp}). The fast path adds the two translations back
explicitly (normalized $\hat{\bm t}_u,\hat{\bm t}_v$, eigenvalue $=$ Rayleigh quotient $\approx10^{-17}$), so the
rigid-mode detection and the output columns are unchanged.

\subsection{Monitoring stability during the run}

The same routine runs inside the loading loop (\code{StabilityMonitor},
\code{src/\allowbreak{}experiments/\allowbreak{}stability\_\allowbreak{}monitor.cpp}), always on relaxed (POST) states:
\begin{itemize}[leftmargin=*]
\item \emph{grid}: every $N$-th load step (\code{-{}-eig-every=N}); at $150\times150$ one computation
costs $\approx0.15$\,s against $\approx0.42$\,s per load step, so $N=5$ adds $\approx7\%$;
\item \emph{refine}: every step while $\lambda_{\min}<0.2\,\lambda_{\mathrm{ref}}$, where
$\lambda_{\mathrm{ref}}$ is the first value computed after the last avalanche
(\code{-{}-eig-refine}); near a saddle-node $\lambda_{\min}\propto\sqrt{\alpha_c-\alpha}$, so the
interesting part of the curve is the last few steps;
\item \emph{before/after avalanche} (\code{-{}-eig-at-avalanche=1}): the relaxed state of the
previous step, POST$(i-1)$, the last stable state before the avalanche, is kept in memory
(positions; its mesh is the one step $i$ starts with), and its eigenpairs are computed
retroactively once step $i$ is recognized as an avalanche; POST$(i)$ is computed as well.
\end{itemize}
Eigenvalues go to \code{eigen\_\allowbreak{}log.csv}; the lowest two nontrivial modes of POST$(i-1)$
go to \code{eigen\_\allowbreak{}modes/\allowbreak{}soft\_\allowbreak{}modes\_XXXXX.vtk}, next to the PRE-avalanche
configuration with the same id. Translations are projected out of the solver and, as a
safeguard, any vector with more than half of its norm in the translations would be skipped.

A first $50\times50$ run (\code{-{}-eig-every=5 -{}-eig-at-avalanche=1}) shows the expected
behaviour, and one instructive case: at steps 1281--1287 $\lambda_{\min}$ decreases from
$3.7\times10^{-3}$ to $1.3\times10^{-3}$ and jumps to $3.4\times10^{-2}$ at step 1288, with a 26\%
stress drop and a 1.4\% energy drop; since remeshing accepted no topological change there,
this instability is \emph{not} counted as an avalanche and no configuration is saved. The
eigenvalue log sees it anyway.

\paragraph{Further options.} If many modes are needed repeatedly, LOBPCG~\cite{knyazev2001}
preconditioned with the PSD-projected Cholesky factor that the L-BFGS solver already has would
avoid the shift search altogether; and since the Cholesky of $K^+$ is computed at every load step
anyway, a single extra factorization of $K+\epsilon I$ (no projection) gives a yes/no stability test
per step for 45\,ms.

%=====================================================================
\section{Related work: what is known and what is new}
\label{sec:literature}
%=====================================================================

It is worth being honest about which parts of this note are textbook material.

\paragraph{Known ingredients.}
\begin{itemize}[leftmargin=*]
\item \emph{L-BFGS with a nontrivial initial inverse Hessian.} L-BFGS~\cite{nocedal1980,liu1989}
only needs a matrix $H_0$ to start its two-loop recursion; that $H_0$ may be any SPD matrix is
standard~\cite{nocedal2006}. Section~\ref{sec:change} (change of variables $\Leftrightarrow$ $H_0=\gamma P^{-1}$) is a
reformulation of this well-known fact, needed only because ALGLIB~\cite{alglib} does not accept a
sparse $H_0$.
\item \emph{The stiffness matrix as $H_0$ in finite elements.} Quasi-Newton methods started from
the factorized stiffness go back to Matthies and Strang~\cite{matthies1979} and are the
``quasi-Newton'' option of commercial FE codes.
\item \emph{Sparse, physics-based preconditioners for L-BFGS in atomistics.} Jiang et
al.~\cite{jiang2004} precondition L-BFGS with part of the Hessian of molecular force fields;
Packwood et al.~\cite{packwood2016} propose a sparse Laplacian-like preconditioner for condensed
matter, show that the gain grows with system size (our $\kappa\sim L^2$ argument), and ship it in
ASE; Mones et al.~\cite{mones2018} build it from the Hessian of the force-field terms, keeping
only the positive definite part, which is what our PSD projection does.
\item \emph{L-BFGS with a prefactorized Laplacian for hyperelasticity}~\cite{liu2017}
(computer graphics): essentially our Laplacian variant.
\item \emph{Element-wise projection of indefinite tangents}~\cite{teran2005} is standard in
``projected Newton'' methods.
\item \emph{Shift-invert Lanczos, inertia, sparse Cholesky}~\cite{ericsson1980,grimes1994,chen2008,amestoy1996}
are classical numerical linear algebra.
\end{itemize}

\paragraph{What seems specific to this work.}
\begin{itemize}[leftmargin=*]
\item The analytic tangent of a $GL(2,\mathbb Z)$-invariant, Lagrange-reduced
energy~\cite{conti2004,baggio2019}, including the piecewise-constant reduction matrix $Z$
(eq.~\eqref{eq:A}), and its use as a preconditioner in a landscape where the Hessian
\emph{jumps} whenever an element changes reduction.
\item The behaviour through avalanches and remeshing: a $K^+$ built at the start of a load step
remains effective across reduction switches, with restarts when it goes stale.
\item The effect of the minimizer on metastable-state selection after instabilities
(Section~\ref{sec:physics}), which matters for avalanche statistics.
\end{itemize}
To our knowledge the Lagrange-reduction plasticity simulations published so far use plain
first-order minimizers (L-BFGS, FIRE-type dynamics~\cite{bitzek2006}). A publication would need
more than a speedup: a size-scaling study ($L=50\ldots500$), comparisons with FIRE, nonlinear CG
and Newton with the same $K$, and statistical equivalence of avalanche distributions over many
seeds; or, for a numerical-analysis journal, a genuinely new ingredient (convergence theory for the
non-smooth reduced energy, instability detection or saddle search built on the factorization,
reuse of the factorization across remeshing, 3D).

%=====================================================================
\section{Implementation map}
%=====================================================================

\begin{center}
\begin{tabular}{p{6.2cm}p{8.3cm}}
\toprule
file / symbol & role\\
\midrule
\code{include/\allowbreak{}optimization/\allowbreak{}PreconditionedLBFGS.h}, \code{src/\allowbreak{}optimization/\allowbreak{}PreconditionedLBFGS.cpp} & everything described in this note\\
\quad\code{element\_\allowbreak{}lagrangian\_\allowbreak{}tangent} & $\mathbb{A}$, eq.~\eqref{eq:A}\\
\quad\code{element\_\allowbreak{}stiffness\_\allowbreak{}from\_\allowbreak{}tangent} & $K^e$, eq.~\eqref{eq:Ke}\\
\quad\code{FastStiffnessAssembler} & cached pattern, parallel assembly of $K^+$ and of the Laplacian\\
\code{include/\allowbreak{}optimization/\allowbreak{}SparseCholesky.h}, \code{src/\allowbreak{}optimization/\allowbreak{}SparseCholesky.cpp} & \code{SparseCholeskyBackend}: CHOLMOD / Eigen factorization, half-solves and full solve\\
\code{include/\allowbreak{}optimization/\allowbreak{}StiffnessSpectrum.h}, \code{src/\allowbreak{}optimization/\allowbreak{}StiffnessSpectrum.cpp} & \code{lowest\_\allowbreak{}stiffness\_\allowbreak{}modes}: lowest eigenpairs of $K$ (Section~\ref{sec:eig})\\
\code{src/\allowbreak{}experiments/\allowbreak{}stability\_\allowbreak{}monitor.cpp} & eigenvalues during the loading, flags \code{-{}-eig-*}\\
\quad\code{transformed\_\allowbreak{}energy} & the ALGLIB callback (change of variables)\\
\quad\code{PreconditionedLBFGS::\allowbreak{}optimize} & setup, ALGLIB run, stopping, restarts\\
\code{BaseLatticeCalculator::\allowbreak{}calculate\_\allowbreak{}dseconderivative\_\allowbreak{}components} & second derivatives without ITensor (+ overrides)\\
\code{energy\_\allowbreak{}functions::\allowbreak{}create\_\allowbreak{}d2phi\_\allowbreak{}components\_\allowbreak{}func} & the six analytic partials\\
\code{relax\_\allowbreak{}configuration()} (\code{common\_\allowbreak{}simulation\_\allowbreak{}helpers.cpp}) & single entry point used by the load steps, \code{memory()} and both remeshing loops; plain L-BFGS unless configured\\
\code{main.cpp} & flags \code{-{}-precond=}, \code{-{}-precond-tol=}, \code{-{}-precond-refresh=}, \code{-{}-precond-from-step=}\\
\code{benchmarks/\allowbreak{}benchmark\_\allowbreak{}preconditioner.cpp} & benchmark, \code{-{}-validate}, \code{-{}-eig N}\\
\code{CMakeLists.txt}, \code{Makefile.cluster} & CHOLMOD detection\\
\code{makefiles\_\allowbreak{}for\_\allowbreak{}cluster/\allowbreak{}job\_\allowbreak{}creator*.py} & preconditioner prompt (default \code{stiffness}, from step 1)\\
\bottomrule
\end{tabular}
\end{center}

\paragraph{Usage.} From the shell:
\begin{lstlisting}[language=bash]
./lattice_triangulation 150 150 positive 42 1 --precond=stiffness --precond-from-step=1
./lattice_triangulation 150 150 positive 42 1      # original plain L-BFGS
\end{lstlisting}

\paragraph{Remeshing.} The remeshing loop calls \code{relax\_\allowbreak{}configuration()} too, so its
relaxations are preconditioned as well, following the same \code{-{}-precond-from-step} rule as the
main solve of that load step. After each remesh the connectivity changes; the solver notices it
(mesh signature) and rebuilds the pattern and the symbolic factorization, which costs a few tens
of milliseconds.

\paragraph{Cluster builds.} Without SuiteSparse the Eigen fallback is used ($\sim8\times$ instead of
$\sim15\times$ at $150\times150$). \code{make -f Makefile.cluster info} tells which one is used.
Because widely included headers changed (a new virtual method, a new member in \code{UserData}),
the first rebuild must start from \code{make clean}; \code{Makefile.cluster} now tracks header
dependencies, \code{Makefile.tchenla} does not yet.

%=====================================================================
\section{Where the remaining time goes, and possible next steps}
%=====================================================================

After this change the solve is no longer the only cost: in the running $150\times150$ simulation a
load step takes $\approx0.45$\,s in total, of which $\approx0.30$\,s is the minimization. The rest is
the energy/stress evaluations before and after, neighbour and defect analysis, remeshing, and
output; before, it was hidden behind 2.9\,s of L-BFGS.

Possible next steps:
\begin{itemize}[leftmargin=*]
\item profile and trim the per-step overhead outside the solver (another factor $\sim1.5$);
\item decide whether instabilities without an accepted remeshing (such as step 1288 above)
should count as avalanches;
\item refresh $K$ during long solves (every $N$ iterations) if avalanche steps ever become the
bottleneck;
\item statistical comparison of avalanche distributions with and without preconditioner over
several seeds, before using the preconditioned runs for publication-level statistics.
\end{itemize}

%=====================================================================
\begin{thebibliography}{99}
\small
\newcommand{\why}[1]{\par\smallskip\noindent\begin{minipage}{\linewidth}\itshape #1\end{minipage}\par\medskip}

\bibitem{nocedal1980}
J.~Nocedal, \emph{Updating quasi-Newton matrices with limited storage},
Math.\ Comp.\ \textbf{35} (1980) 773--782.
\why{The original limited-memory BFGS update and the two-loop recursion. The recursion takes an
arbitrary initial matrix $H_0$; everything in Section~\ref{sec:change} relies on this.}

\bibitem{liu1989}
D.~C.~Liu and J.~Nocedal, \emph{On the limited memory BFGS method for large scale optimization},
Math.\ Programming \textbf{45} (1989) 503--528.
\why{Convergence theory and the scaling $H_0=\gamma_k I$, $\gamma_k=\s\T\y/\y\T\y$, that ALGLIB uses
by default and that our change of variables turns into $H_0=\gamma_k P^{-1}$.}

\bibitem{nocedal2006}
J.~Nocedal and S.~J.~Wright, \emph{Numerical Optimization}, 2nd ed., Springer (2006), Chs.~6--7.
\why{Textbook treatment of (L-)BFGS, preconditioning and its equivalence with a linear change of
variables, and of Newton and trust-region methods, the natural competitors once $K$ is available.}

\bibitem{alglib}
S.~Bochkanov, \emph{ALGLIB numerical library}, \url{https://www.alglib.net}, manual of
\code{minlbfgs} (version~3.19).
\why{The optimizer we use. Its preconditioners (default, diagonal, dense Cholesky, scale) do not
include a sparse factorization, which is why we precondition from outside
(Section~\ref{sec:change}).}

\bibitem{matthies1979}
H.~Matthies and G.~Strang, \emph{The solution of nonlinear finite element equations},
Int.\ J.\ Numer.\ Meth.\ Engng.\ \textbf{14} (1979) 1613--1626.
\why{BFGS started from the factorized FE stiffness matrix, updated with low-rank corrections instead
of refactorizing: the direct ancestor of what we do. The idea is standard in structural mechanics
codes; our contribution is its application to a Lagrange-reduced, non-smooth energy.}

\bibitem{jiang2004}
L.~Jiang, R.~H.~Byrd, E.~Eskow and R.~B.~Schnabel, \emph{A preconditioned L-BFGS algorithm with
application to molecular energy minimization}, Tech.\ Rep.\ CU-CS-982-04, Univ.\ of Colorado (2004).
\why{Preconditioned L-BFGS with a sparse preconditioner built from parts of the (partially
separable) molecular energy; speedups 2--7 in CPU time for proteins. Same philosophy as ours: use
the structure of the Hessian, not a generic diagonal.}

\bibitem{packwood2016}
D.~Packwood, J.~Kermode, L.~Mones, N.~Bernstein, J.~Woolley, N.~Gould, C.~Ortner and G.~Cs\'anyi,
\emph{A universal preconditioner for simulating condensed phase materials},
J.\ Chem.\ Phys.\ \textbf{144} (2016) 164109.
\why{The closest reference for the size-scaling argument of Section~\ref{sec:slow}: for materials the condition
number of the Hessian grows with system size, and a sparse Laplacian-type preconditioner built
from the neighbour graph removes this growth. Implemented in ASE for L-BFGS and FIRE. Our
\code{laplacian} option is the FE analogue; our \code{stiffness} option uses the exact Hessian.}

\bibitem{mones2018}
L.~Mones, C.~Ortner and G.~Cs\'anyi, \emph{Preconditioners for the geometry optimisation and
saddle point search of molecular systems}, Sci.\ Rep.\ \textbf{8} (2018) 13991.
\why{Sparse preconditioners built from the Hessian of force-field terms, keeping only the positive
definite part: conceptually the same as our element-wise PSD projection (Section~\ref{sec:psd}). They also use
the preconditioner for saddle searches, a possible extension for avalanches.}

\bibitem{liu2017}
T.~Liu, S.~Bouaziz and L.~Kavan, \emph{Quasi-Newton methods for real-time simulation of
hyperelastic materials}, ACM Trans.\ Graph.\ \textbf{36} (2017) 23.
\why{L-BFGS for hyperelastic FE with a prefactorized, constant Laplacian-type matrix as initial
Hessian; shows the large gains of a good $H_0$ in elasticity. Essentially our \code{laplacian}
variant; our \code{stiffness} variant goes further by refactorizing the true tangent every step.}

\bibitem{teran2005}
J.~Teran, E.~Sifakis, G.~Irving and R.~Fedkiw, \emph{Robust quasistatic finite elements and flesh
simulation}, Proc.\ ACM SIGGRAPH/Eurographics Symp.\ Computer Animation (2005) 181--190.
\why{Makes indefinite element tangents definite element by element (diagonalization and clamping),
so that Newton steps are descent directions. Our PSD projection of the $4\times4$ tangent
$\mathbb A$ is the same idea.}

\bibitem{bitzek2006}
E.~Bitzek, P.~Koskinen, F.~G\"ahler, M.~Moseler and P.~Gumbsch, \emph{Structural relaxation made
simple}, Phys.\ Rev.\ Lett.\ \textbf{97} (2006) 170201.
\why{FIRE, the damped-dynamics minimizer widely used in athermal quasi-static simulations of
plasticity and amorphous solids; the obvious baseline for any comparison beyond plain L-BFGS.}

\bibitem{ericsson1980}
T.~Ericsson and A.~Ruhe, \emph{The spectral transformation Lanczos method for the numerical
solution of large sparse generalized symmetric eigenvalue problems}, Math.\ Comp.\ \textbf{35}
(1980) 1251--1268.
\why{Shift-invert (spectral transformation) Lanczos: eigenvalues near a shift $\sigma$ become the
dominant ones of $(K-\sigma I)^{-1}$. Basis of Section~\ref{sec:eig}.}

\bibitem{grimes1994}
R.~G.~Grimes, J.~G.~Lewis and H.~D.~Simon, \emph{A shifted block Lanczos algorithm for solving
sparse symmetric generalized eigenproblems}, SIAM J.\ Matrix Anal.\ Appl.\ \textbf{15} (1994)
228--272.
\why{The ``industrial'' shift-invert Lanczos of FE codes; uses the inertia of the $LDL\T$
factorization of $K-\sigma I$ (Sylvester's law) to count eigenvalues below a shift and to steer
the shifts. Our shift search is the simplest version of this: a Cholesky failure means an
eigenvalue below $\sigma$.}

\bibitem{knyazev2001}
A.~V.~Knyazev, \emph{Toward the optimal preconditioned eigensolver: locally optimal block
preconditioned conjugate gradient method}, SIAM J.\ Sci.\ Comput.\ \textbf{23} (2001) 517--541.
\why{LOBPCG: computes the lowest eigenpairs with a preconditioner instead of exact solves. A
natural option for us, since the PSD-projected Cholesky factor is already available.}

\bibitem{amestoy1996}
P.~R.~Amestoy, T.~A.~Davis and I.~S.~Duff, \emph{An approximate minimum degree ordering
algorithm}, SIAM J.\ Matrix Anal.\ Appl.\ \textbf{17} (1996) 886--905.
\why{AMD, the fill-reducing ordering $\Pi$ used before the Cholesky factorization
(Section~\ref{sec:cholesky}).}

\bibitem{chen2008}
Y.~Chen, T.~A.~Davis, W.~W.~Hager and S.~Rajamanickam, \emph{Algorithm 887: CHOLMOD, supernodal
sparse Cholesky factorization and update/downdate}, ACM Trans.\ Math.\ Softw.\ \textbf{35} (2008) 22.
\why{The supernodal sparse Cholesky we use (SuiteSparse). Its update/downdate feature could
later allow modifying the factor after local changes instead of refactorizing.}

\bibitem{spectra}
Y.~Qiu, \emph{Spectra: C++ library for large scale eigenvalue problems},
\url{https://spectralib.org}.
\why{Header-only implicitly restarted Lanczos/Arnoldi (ARPACK-like) used by the eigenvalue code;
we give it our own Cholesky-based shift-invert operator.}

\bibitem{conti2004}
S.~Conti and G.~Zanzotto, \emph{A variational model for reconstructive phase transformations in
crystals, and their relation to dislocations and plasticity}, Arch.\ Rational Mech.\ Anal.\
\textbf{173} (2004) 69--88.
\why{The $GL(2,\mathbb Z)$-invariant energy and the Lagrange reduction used in this code.}

\bibitem{baggio2019}
R.~Baggio, E.~Arbib, P.~Biscari, S.~Conti, L.~Truskinovsky, G.~Zanzotto and O.~U.~Salman,
\emph{Landau-type theory of planar crystal plasticity}, Phys.\ Rev.\ Lett.\ \textbf{123} (2019)
205501.
\why{The Landau-type model of crystal plasticity (Lagrange-reduced energy, quasi-static loading,
avalanches) that the simulations in this note implement; the physics context of the solver.}

\end{thebibliography}

\end{document}
